\section{Introduction}\label{sec:intro}
Residents in Taichung often bring household waste directly to a collection vehicle at a designated location or along a route. The city's collection information service covers 29 districts, more than 370 routes, and over 20,000 collection points, according to its 2026 announcement \cite{tcepb2026cleaner}. Collection windows can be short. During the 2026 Lunar New Year, the city described stops of roughly 5--30 minutes at fixed collection locations \cite{tcepb2026spring}. Missing a vehicle can therefore require a resident to wait for another service or seek another point.

The official service provides collection information and vehicle tracking, while a third-party app provides nearby locations and walking times \cite{tcepb2026cleaner,appstore_private}. The project's review found that users still need to compare those pieces of information to decide where to go. A location with the shortest walk can be a poor choice if its vehicle is about to arrive or has already passed. This makes the decision dependent on time as well as location.

Public data create opportunities for such a service. Taichung publishes vehicle observations and a separate dataset of collection locations and schedules \cite{dataset83558,dataset84004}. These datasets have practical constraints: a position may be old, a scheduled time may differ from actual arrival, and a textual address must be converted into coordinates before spatial routing. Connecting these sources requires attention to their timestamps and identifiers.

GarbageGo addresses this scenario through a planned application that ranks collection points according to estimated time feasibility. It combines the user's location, a walking route, an estimated truck arrival, and a buffer for uncertainty. The intended result includes a reason for the recommendation, the age or source of its data, and alternatives when available. The following sections compare related work, describe the design, and examine the project's existing data audit.
